\subsection*{A certified minimax recoverability floor}
\label{sec:floor}
Proposition~\ref{thm:perlead} says \emph{whether} a target lead's dipolar coordinate is recoverable
($\eta{=}0$) and \emph{how well-conditioned} the estimate is ($\kappa$). We now turn the binary
verdict into a \emph{quantitative} per-lead risk floor in physical units that lower-bounds
\emph{every} reconstructor, linear or learned. Fix segment $s$ and observed set $S$; write the
dipole prior second moment $\bm\Sigma_d=\mathbb E[\bm d\bm d^\top]$, the observed-dipole row-space
projector $\bm P=\bm M_{s,S}^{+}\bm M_{s,S}$, and $\bm Q=\bm I-\bm P$. The target is the dipolar
functional of lead $\ell$, $\theta_\ell=\bm e_\ell^\top\bm M_s\bm d$, and the prior-conditional
ambiguity of Sec.~\ref{sec:cert} is
\[
a_{s,\ell}(S)^2 \;=\; \bm e_\ell^\top\bm M_s\,\bm\Sigma_{Q\mid P}\,\bm M_s^\top\bm e_\ell,
\qquad
\bm\Sigma_{Q\mid P}=\bm Q\bm\Sigma_d\bm Q-\bm Q\bm\Sigma_d\bm P(\bm P\bm\Sigma_d\bm P)^{+}\bm P\bm\Sigma_d\bm Q .
\]

\begin{theoremc}[Minimax per-lead recoverability floor]\label{thm:floor}
Let $\mathcal P(\bm\Sigma_d)=\{\pi:\ \mathbb E_\pi[\bm d\bm d^\top]\preceq\bm\Sigma_d\}$ be the class of
priors with second moment dominated by $\bm\Sigma_d$.
\begin{enumerate}[label=(\roman*),leftmargin=1.6em]
\item \textbf{Noise-free dipolar floor (nonlinear estimators).} For \emph{any} estimator
$\hat\theta_\ell(\bm y_S)$ (measurable, possibly nonlinear/learned) of $\theta_\ell$ from the exact
observation $\bm y_S=\bm M_{s,S}\bm d$,
\[
\inf_{\hat\theta_\ell}\ \sup_{\pi\in\mathcal P(\bm\Sigma_d)}\ \mathbb E_\pi\big[(\hat\theta_\ell-\theta_\ell)^2\big]
\;=\; a_{s,\ell}(S)^2 ,
\]
attained by the least-favourable Gaussian prior $\mathcal N(\bm 0,\bm\Sigma_d)$ and its linear
posterior-mean estimator. In particular $a_{s,\ell}(S)^2=0$ iff $\eta_{s,\ell}(S)=0$, and every
reconstructor incurs risk $\ge a_{s,\ell}(S)^2$ under the least-favourable prior in the class.
\item \textbf{Noise/residual inflation.} With observed noise and non-dipolar residual of covariance
$\bm N$ on $S$, the risk floor lifts to
$a_{s,\ell}(S)^2+\lambda_{\min}$-scale term bounded through
$\kappa_{s,\ell}(S)=\|\bm e_\ell^\top\bm M_s\bm M_{s,S}^{+}\|_2$: the amplified error is
$\bm e_\ell^\top\bm M_s\bm M_{s,S}^{+}(\bm{\mathrm{Sel}}_S\bm r_s+\bm n)$, of variance
$\le\kappa_{s,\ell}(S)^2\,\|\bm N\|_2$ (Prop.~\ref{thm:perlead}(ii)).
\item \textbf{Finite-sample subspace floor.} With $\bm M_s$ replaced by its rank-3 PCA estimate
$\widehat{\bm M}_s$ from $N$ records,
$\big|a_{s,\ell}(S;\widehat{\bm M}_s)^2-a_{s,\ell}(S;\bm M_s)^2\big|\le C\,\sin\Theta(\widehat{\bm M}_s,\bm M_s)\,\|\bm\Sigma_d\|_2$,
where $\sin\Theta$ is the principal-angle perturbation (Davis--Kahan~\cite{davis1970rotation}). The
conservative floor $a^{-}_{s,\ell}(S)=\min$ over the record-bootstrap confidence set of
$\widehat{\bm M}_s$ then lower-bounds the risk uniformly in the estimated subspace.
\end{enumerate}
\end{theoremc}

\noindent\emph{Proof sketch.} (i) Gaussian conditioning gives
$\mathrm{Var}(\theta_\ell\mid\bm P\bm d)=\bm e_\ell^\top\bm M_s\bm\Sigma_{Q\mid P}\bm M_s^\top\bm e_\ell$
with the posterior mean Bayes-optimal under squared loss, so the Bayes risk equals $a_{s,\ell}^2$;
estimating the \emph{linear} functional $\theta_\ell=\bm e_\ell^\top\bm M_s\bm d$ under a
convex, quadratically-constrained parameter class makes the nonlinear minimax risk equal the linear
one, with a least-favourable Gaussian prior (Donoho--Liu~\cite{donohoLiu1991}; Donoho, optimal
recovery~\cite{donoho1994}), giving the sup--inf identity and the bound for arbitrary $\hat\theta_\ell$.
(ii) Cauchy--Schwarz/$\sigma_{\min}$ on the amplified-error term, as in
Prop.~\ref{thm:perlead}(ii). (iii) Weyl/Davis--Kahan perturbation of the projectors and Schur
complement as functions of $\widehat{\bm M}_s$, propagated through the quadratic form.\hfill$\square$

\smallskip
\noindent\textbf{What is and is not certified.} The floor is \emph{minimax over the moment class}
$\mathcal P(\bm\Sigma_d)$, not distribution-free: a reconstructor exploiting higher moments the true
(non-Gaussian) $\bm d$ happens to carry could beat $a_{s,\ell}^2$ under the \emph{true} prior, not
the least-favourable one. The certified quantity is the \emph{dipolar} functional $\theta_\ell$,
where the algebra is exact; a floor on the \emph{full} lead $\bm e_\ell^\top\bm L_s$ would require
bounding $\operatorname{tr}\mathrm{Cov}(\bm r_\ell\mid\bm y_S)$, which we decline to certify and
instead meter empirically with the distribution-free intervals of Sec.~\ref{sec:calib}. So
``no estimator beats $a_{s,\ell}(S)$ over the moment class'' is a target-specific, finite-sample
statement about the identifiable dipolar component; we verify it numerically under the
least-favourable Gaussian prior (\texttt{tests/test\_floor.py}: zero floor violations). On real ECG
the floor is \emph{operational and tight}: the theoretical ambiguity $a_{s,\ell}$ strongly predicts
the measured per-lead dipolar error (Spearman $\CertValSpearAmb$, Pearson $\CertValPearAmb$;
Fig.~\ref{fig:certval}), the identifiable/unidentifiable split is separated with AUC
$\CertValAUC$, and the pure-dipolar reconstructor is bounded by the floor while data-driven
reconstructors sit right at it (median gap $\CertValFloorGap$\,mV)---beating the moment-class worst
case only marginally by exploiting non-Gaussian population structure, exactly the residual the
calibrated intervals of Sec.~\ref{sec:calib} meter. This turns the empirically similar total error
across reconstructors (Sec.~\ref{sec:exp}) into a target-specific certified minimax floor.
Crucially the floor also binds a \emph{learned} reconstructor: a trained arbitrary-mask diffusion
model (DDPM, honest guidance $w{=}1$), evaluated against the \emph{same} floor on held-out records,
sits above it in every one of $\FloorDiffNcells$ underdetermined cells (floor-violation fraction
$\FloorDiffViol$, median gap $\FloorDiffGap$\,mV), and the certificate's identifiable/unidentifiable
split predicts its error---median dipolar-projection RMSE $\FloorDiffEtaZero$\,mV where
$\eta{=}0$ versus $\FloorDiffEtaPos$\,mV where $\eta{>}0$ (\texttt{experiments/certificate\_floor\_diffusion.py}).
Notably the generative model does \emph{not} beat the moment-class floor at all, whereas the linear
ridge does on a minority of cells---the worst-case bound is loose only for the estimator explicitly
exploiting the non-Gaussian structure it excludes.
